\documentclass{aastex631}

\usepackage{amsmath,amssymb}
\usepackage{natbib}
\usepackage{booktabs}

\newcommand{\dd}{\mathrm{d}}
\newcommand{\pp}{\partial}
\newcommand{\Erad}{E_{\rm rad}}
\newcommand{\Eg}{E_{\rm g}}
\newcommand{\Elost}{E_{\rm lost}}
\newcommand{\Tint}{T_{\rm int}}
\newcommand{\Teff}{T_{\rm eff}}
\newcommand{\epsA}{\varepsilon_A}
\newcommand{\epsS}{\varepsilon_S}
\newcommand{\dE}{\delta E}

\shorttitle{ORCHARD Energy Conservation}
\shortauthors{Tejada Arevalo}

\begin{document}

\title{Energy Conservation Accounting and Error Budget in the ORCHARD Planetary Evolution Code}

\author{Roberto Tejada Arevalo}

\begin{abstract}
We describe the global energy-conservation accounting of the ORCHARD
planetary evolution code: the definitions of the energy ledger, a set of
accounting corrections applied in 2026 July, and an exact decomposition of
the conservation residual.  The residual splits into three independently
measurable defect channels: an equation-of-state (EOS) table thermodynamic
inconsistency $\epsA$, a transport-solver entropy-closure defect $\epsS$,
and a work-versus-gravity discretization defect $C$ of the hydrostatic
scheme.  The decomposition closes to machine precision by construction and
cleanly separates physics-input error from numerical error.  Applied to
homogeneous and inhomogeneous Jupiter models on a resolution ladder
($N=300$--$1200$ mass zones), it shows that $\epsA$ is
resolution-independent ($\simeq +9.3\times10^{-3}$ per unit radiated
energy), $\epsS$ is small ($\simeq +1.6\times10^{-3}$), and $C$ scales as
$-4.7/N$.  The two dominant
terms have opposite signs and comparable magnitudes at customary
resolutions.  The raw residual is therefore misleadingly small near
$N\simeq450$, through accidental cancellation.  We trace the $O(1/N)$ term
to a single first-order (one-sided) pressure factor in the gravity term of
the Henyey difference equations.  Symmetrizing it to the midpoint form
reduces $|C|$ by a factor of $\sim$200, renders the residual
resolution-independent at the EOS floor, and accelerates the convergence
of observables such that $N=300$ midpoint models match
Richardson-extrapolated original-scheme results.  The midpoint scheme is
presented here as a validated experiment.  It is not yet the default.
\end{abstract}

\section{The Energy Ledger}
\label{sec:ledger}

ORCHARD evolves the entropy $S$, composition $(Y,Z)$, and hydrostatic
structure $(P,r)$ on a Lagrangian mass grid.  Total energy is not a solved
variable.  Global energy conservation is therefore a \emph{diagnostic}: an
independent audit of the solution, evaluated from the state at every
accepted timestep.

The ledger tracks three quantities:
\begin{align}
U(t)      &= \sum_k u_k\,\dd m_k, &
\Eg(t)    &= \sum_k \frac{G m_k\,\dd m_k}{r_k}, &
\Elost(t) &= \sum_{\rm steps} L_{\rm surf}\,\Delta t,
\end{align}
where $u_k$ is the specific internal energy from the EOS at $(P_k, T_k,
Y_k, Z_k)$, $m_k$ and $r_k$ are cell-centered enclosed mass and radius,
and $L_{\rm surf}$ is the surface luminosity removed by the transport
step.  With the sign convention that $\Eg>0$, the total energy is
$E = U - \Eg$ and the conservation residual is
\begin{equation}
\dE(t) \;=\; \bigl[\,U(t) - \Eg(t) + \Elost(t)\,\bigr] \;-\;
             \bigl[\,U(0) - \Eg(0)\,\bigr].
\label{eq:residual}
\end{equation}
Perfect conservation gives $\dE = 0$ at all times: the energy radiated
equals the energy the interior released.

\subsection{Normalization}
\label{sec:normalization}

Two normalizations of the same residual appear in the code output:
\begin{equation}
\mathrm{d}E_G \equiv \frac{\dE}{e_0}, \qquad
\frac{\dE}{\Erad} \equiv \frac{\dE}{\Elost(t)},
\end{equation}
with $e_0 = U(0)-\Eg(0)$.  The historical metric $\mathrm{d}E_G$ is
\emph{deprecated for physical interpretation}.  Because $u$ is defined
only up to an additive constant per EOS table, $e_0$ (and hence
$\mathrm{d}E_G$) depends on the table's zero point.  Between the
homogeneous and inhomogeneous Jupiter reference models, $e_0$ differs by a
factor of 1.54 while the physical error is nearly identical --- the
apparent factor-of-4 difference in their $\mathrm{d}E_G$ values that
originally motivated this work was entirely a normalization artifact.
Because
$e_0<0$ for a bound planet, $\mathrm{d}E_G$ also has the opposite sign of
the residual.

$\dE/\Erad$ is zero-point-free and grades the calculation by the energy it
actually processed; empirically, all defect channels accrue in proportion
to $\Erad(t)$, making $\dE/\Erad$ approximately stationary in time and
comparable across planets.  It is the quantity quoted throughout.

\section{Accounting Corrections (2026 July)}
\label{sec:fixes}

Four corrections were applied to make the ledger internally consistent
(commit \texttt{a584ac5}).  All are bookkeeping-only and were verified to
leave the evolved physics bit-identical.

\paragraph{(i) Radiated-energy quadrature.}
The transport solver is backward-Euler: it removes the surface flux of the
\emph{new} state over the whole step.  The consistent quadrature is the
rectangle
\begin{equation}
\Delta \Elost = L_{\rm surf}^{\,n+1}\,\Delta t .
\end{equation}
The legacy ledger instead used a trapezoid of the \emph{average} of two
different luminosity channels --- the post-hydrostatic
$4\pi r_0^2 \sigma \Tint^4$ (evaluated at a cell-center radius, after the
atmospheric boundary update) and the transport-solver surface flux.  That
mixture meters a quantity that is not the heat actually drained from the
entropy field.  The measured bias is $O(10^{-4})\,\Erad$, model-dependent
in sign.  The historical channels are retained as debug diagnostics.

\paragraph{(ii) Regrid mass refresh.}
The cell-centered enclosed-mass array entering $\Eg$ was built once at
startup and not refreshed when adaptive regridding rebound the mass grid;
$\Eg$ then mixed stale masses with new radii.  (Latent for all
non-regridding models.)

\paragraph{(iii) Restart fail-fast.}
Restarting from saved output cannot reconstruct the ledger histories, and
previously crashed with an \texttt{IndexError} at the first accepted step.
Consequently, \texttt{restart\_flag=True} now raises immediately.  ORCHARD
policy is to always run models fresh.

\paragraph{(iv) Reporting.}
The final summary reports $\dE/\Erad$ first, with $\mathrm{d}E_G$ labeled
as zero-point-dependent.

\section{Exact Error Decomposition}
\label{sec:decomposition}

The residual (Eq.~\ref{eq:residual}) is the sum of the violations of
three physical identities, one per code subsystem.  Writing
$W = \sum_{\rm steps}\sum_k (P/\rho^2)_k\,\Delta\rho_k\,\dd m_k$ for the
cumulative compression work and $\mu$-terms for the composition
(chemical-potential) contributions:

\begin{enumerate}
\item \textbf{EOS first law} (violated by $\epsA$).  A thermodynamically
consistent table satisfies
$\dd u = T\,\dd s + (P/\rho^2)\,\dd\rho + \mu_Y \dd Y + \mu_Z \dd Z$
exactly, so
\begin{equation}
\Delta U = \!\int\!\! T\,\dd S\,\dd m + W + \mu\text{-terms} + \epsA .
\end{equation}
$\epsA$ is the Maxwell-relation defect of the table: the amount by which
the tabulated $u$ fails to be the potential of the tabulated $(T, P,
\rho, s)$ along the evolutionary path.  It vanishes for a perfect table
at any resolution and timestep.

\item \textbf{Transport heat balance} (violated by $\epsS$).  The entropy
equation drains heat only through the surface:
\begin{equation}
\int T\,\dd S\,\dd m + \mu\text{-terms} = -\Erad + \epsS .
\end{equation}
$\epsS$ is the closure defect of the implicit transport step; see
Section~\ref{sec:epsS}.

\item \textbf{Structure work--energy identity} (violated by $C$).  For
quasi-static contraction in hydrostatic equilibrium, integration by parts
of the HSE equation gives $\Delta(-\Eg) = -W$ up to a negligible surface
term ($P_{\rm atm}\,\Delta V \sim 10^{-8}\,\Erad$), so
\begin{equation}
\Delta(-\Eg) = -W + C .
\end{equation}
$C$ measures the failure of the \emph{discrete} state sequence to satisfy
the discrete analog of this identity.
\end{enumerate}

Adding the three identities, the $T\dd S$ and $W$ terms cancel and
\begin{equation}
\boxed{\;\dE \;=\; \epsA + \epsS + C\;}
\label{eq:budget}
\end{equation}
exactly.  In practice $\epsS$ is read from the code's own per-timestep
accumulation (the \texttt{dE\_S} diagnostic:
$\epsS = \mathrm{d}E_S \times \Elost$, float64, chemical terms included);
$W$ is reconstructed from saved profiles with save-interval midpoints;
$C$ follows from the stored $\Eg$; and $\epsA$ is the remainder.  The
identity closes to $\sim$$10^{-16}$ by construction.  The $\epsA$/$C$
\emph{split} shares the quadrature error of $W$, and is validated by the
resolution-independence of $\epsA$ below.

\section{Results on the Resolution Ladder}
\label{sec:results}

Homogeneous (no-rain, $10\,M_\oplus$ core) and inhomogeneous
(Gaussian-$Z$, He-rain) Jupiter models were evolved to 5\,Gyr on a ladder
$N = 300$--$1200$.  Table~\ref{tab:ladder} lists the final budget terms.

\begin{deluxetable}{rrrrr}
\tablecaption{Energy budget vs.\ resolution (homogeneous Jupiter, original
scheme, $\times10^{-3}\,\Erad$).\label{tab:ladder}}
\tablehead{\colhead{$N$} & \colhead{$\dE/\Erad$} & \colhead{$\epsA$} &
\colhead{$\epsS$} & \colhead{$C$}}
\startdata
 300 & $-4.90$ & $+9.26$ & $+1.63$ & $-15.79$ \\
 450 & $+0.60$ & $+9.27$ & $+1.62$ & $-10.29$ \\
 600 & $+3.35$ & $+9.26$ & $+1.70$ & $-7.61$  \\
 900 & $+5.86$ & $+9.26$ & $+1.68$ & $-5.07$  \\
1200 & $+7.15$ & $+9.25$ & $+1.63$ & $-3.74$  \\
\enddata
\tablecomments{The inhomogeneous model gives $\epsA+\epsS = +11.7$ and
$+11.5$ at $N=300$ and 600, with $C=-16.0$ and $-8.3$: the same scaling
and a slightly higher, but equally $N$-independent, EOS floor.}
\end{deluxetable}

Three facts follow.  (1)~$\epsA$ is constant to $0.3\%$ over a factor of
4 in $N$ --- it is a property of the EOS data files, not of the numerics.
(2)~$C \simeq +0.3\times10^{-3} - 4.7/N$: an almost pure first-order
spatial error.  (3)~The signed residual rises monotonically along
$\mathrm{floor} + \mathrm{const}/N$ toward the EOS floor
$\epsA+\epsS \simeq +10.9\times10^{-3}$, crossing zero near
$N\!\approx\!450$.

\subsection{The cancellation trap}
The near-zero residual at $N\!\approx\!450$ is \emph{accidental}.  Two
independent error sources of opposite sign and comparable magnitude
cancel.  There is no mechanism linking them.  The sign of $\epsA$ is
set by the direction of the table's Maxwell violation along cooling
paths, while the sign of $C$ is set by the one-sidedness of the discrete
gravity term (Section~\ref{sec:rootcause}).  Historically, quoting the
raw residual at customary resolutions ($N=300$--600) therefore
\emph{understated} both errors.  Conversely, refining the grid appeared to
``worsen'' conservation as the cancellation was progressively removed.
Tuning $N$ to minimize $|\dE|$ optimizes for cancellation, not accuracy.

\subsection{Spatial structure of $C$}
A per-cell decomposition
$c_k = G m_k \dd m_k\,\Delta(1/r_k) + \sum_{\rm steps}(P/\rho^2)_k
\Delta\rho_k\,\dd m_k$
shows that the defect \emph{density} is a smooth, planet-wide dipole ---
positive through the deep envelope (peaking just above the core
boundary), crossing zero near $m/M\simeq0.55$, negative in the outer
envelope --- whose shape is $N$-independent and whose cumulative sum
sweeps to $\pm0.3\,\Erad$ before cancelling.  This dipole is physical
(local gravitational energy release need not equal local compression
work; energy is transported between shells, and only the global integral
vanishes), and $C$ is therefore the $O(h)$ error of discretely summing a
large cancelling integrand.  Consequently no local mesh repair can remove
it, and (as verified with five quadrature variants agreeing to four
significant figures, plus a boundary-centered variant matching to three
digits) no re-quadrature of $\Eg$ can either: the leak is in the discrete
states themselves.

\section{Root Cause of $C$ and the Midpoint-Gravity Scheme}
\label{sec:rootcause}

The Henyey pressure equation across interface $k$ is, in the code's
log variables $(x=\ln r_{\rm b},\ y=\ln P)$ and neglecting rotation,
\begin{equation}
y_{k-1} - y_k + \frac{G}{8\pi}\,
\bigl(m_{{\rm b},k-1}-m_{{\rm b},k+1}\bigr)\, m_{{\rm b},k}\,
e^{-y_k - 4x_k} = 0 .
\label{eq:hse}
\end{equation}
The gravity term carries the \emph{one-sided} pressure factor $e^{-y_k}$
(the inner cell only), whereas the rotation term in the same equation
already uses the midpoint factor $e^{-(y_{k-1}+y_k)/2}$.  A one-sided
factor is first-order.  That is precisely the signature of the observed
clean $1/N$ energy defect.

Symmetrizing the gravity factor to $e^{-(y_{k-1}+y_k)/2 - 4x_k}$ (with
the three consistent Jacobian entries) and rerunning the homogeneous
model gives Table~\ref{tab:midpoint}.

\begin{deluxetable}{lrrrr}
\tablecaption{Midpoint-gravity experiment (homogeneous Jupiter,
$\times10^{-3}\,\Erad$; radius in $R_{\rm J}$ at 5\,Gyr).
\label{tab:midpoint}}
\tablehead{\colhead{Scheme} & \colhead{$N$} & \colhead{$C$} &
\colhead{$\dE/\Erad$} & \colhead{$R(5\,{\rm Gyr})$}}
\startdata
original & 300  & $-15.79$ & $-4.90$  & 1.012 \\
original & 600  & $-7.61$  & $+3.35$  & 1.005 \\
original & 1200 & $-3.74$  & $+7.15$  & 1.001 \\
midpoint & 300  & $-0.08$  & $+10.90$ & 0.997 \\
midpoint & 600  & $+0.07$  & $+11.03$ & 0.998 \\
\enddata
\end{deluxetable}

The effects are threefold.  (1)~$|C|$ drops by a factor of $\sim$200, to
a level with no remaining $N$-scaling.  Its sign flips between $N=300$ and
600, at the save-cadence quadrature floor of the diagnostic itself.
(2)~The residual becomes resolution-independent at the EOS floor,
$+10.9\times10^{-3}$, at every age.  The time-resolved
$\dE/\Erad(t)$ curves for $N=300$ and $N=600$ overlie each other and the
floor from the first save onward.  (3)~The scheme is genuinely more
accurate, not merely better-bookkept: Richardson extrapolation of the
original radius ladder gives $R_\infty \simeq 0.997\,R_{\rm J}$, which
the midpoint scheme reaches already at $N=300$; the original scheme at
$N=1200$ still carries a $+0.004\,R_{\rm J}$ error.  $\Teff$ shifts by
$+1.3$\,K at $N=300$.  Henyey iteration counts and wall time are
unchanged.

\emph{Status:} the midpoint form changes all results at the $\sim$1\%
radius level at $N=300$ and therefore breaks comparability with prior
campaign grids.  It is implemented and validated only as an isolated
experiment.  The recommended adoption path is a configuration flag
(\texttt{hse\_gravity\_centering = legacy | midpoint}, default
\texttt{legacy}), with a full-suite validation (inhomogeneous, rotating +
theory-of-figures, sub-Neptune, super-Earth models) before the default is
flipped in a public release.

\section{The Transport Closure Defect $\epsS$}
\label{sec:epsS}

$\epsS \simeq +1.6\times10^{-3}\,\Erad$ is the smallest budget term.  It
is also the only one attributable to the transport subsystem.  An
instrumented
run accumulating the heat with three temperature weightings settles its
origin completely:
\begin{equation}
\epsS(T_{\rm old}) = +6.8\times10^{-8}, \quad
\epsS(T_{\rm mid}) = +1.627\times10^{-3}, \quad
\epsS(T_{\rm new}) = +3.254\times10^{-3},
\end{equation}
with $\epsS(T_{\rm new}) = 2\,\epsS(T_{\rm mid})$ exactly.  The
old-state weighting closes the identity to the Newton-tolerance floor:
the solver's discrete entropy equation carries its $1/T$ factor at the
\emph{beginning-of-step} temperature, and its own heat ledger
($\sum T^{\,n}\Delta S\,\dd m$ against the surface flux) telescopes
perfectly.  The entire $\epsS$ is therefore the first-order
temperature-weighting offset between the solver's $T^{\,n}$ extraction
and the second-order midpoint path heat $\sum T_{\rm mid}\Delta S\,\dd
m$.  This includes the part a coarser cap-halving test had suggested was
$\Delta t$-independent: the caps bind only late, and the early
error-controlled phase, identical in both runs, dominates the
accumulation.  During cooling $T^{\,n} > T_{\rm mid}$, so the scheme
systematically over-extracts heat by $0.16\%$ of $\Erad$.

The exact analog of the gravity-centering fix is implemented behind the
opt-in flag \texttt{[transport] time\_centered\_heat\_factor} (default
\texttt{False}, bit-identical legacy behavior): the heat terms of the
entropy residual (flux divergence and composition $\dd U$ terms) are
metered at a time-centered temperature while the fluxes remain fully
implicit.  Because the flux treatment --- the stiff part --- stays
backward-Euler, this does \emph{not} incur the stability risk of a
Crank--Nicolson flux scheme.  Only the heat-metering weighting is
centered.

The centering estimate matters.  A first implementation used the
solver's internal constant-structure linearization
$\Delta T \simeq (T/C_V)\,\Delta S$ and \emph{overshot}.  That form
overestimates the realized per-step $|\Delta T|$ by a factor
$\sim$2.7 (the true step is closer to constant pressure and is partially
offset by compression heating), driving $\epsS$ from
$+1.63\times10^{-3}$ to $-2.7\times10^{-3}$.  The adopted
implementation instead predicts the midpoint from the \emph{previous
accepted step's actual temperature rate},
$T_{1/2} = T^{\,n} + \tfrac{1}{2}\Delta t\,(\dd T/\dd t)_{\rm prev}$ ---
second-order for smooth cooling and independent of the Newton unknowns,
so that the analytic Jacobian needs no modification.  It falls back to
the legacy weighting on the first step and immediately after a regrid.

Measured on the homogeneous $N=300$ Jupiter: $\epsS$ drops from
$+1.63\times10^{-3}$ to $-4.4\times10^{-5}$ (a factor of $\sim$37, at
the predictor's curvature floor).  With both centering flags active the
full budget becomes
$\dE/\Erad = +9.18\times10^{-3}$ against $\epsA = +9.31\times10^{-3}$,
with $C = -7.7\times10^{-5}$ and $\epsS = -4.8\times10^{-5}$: the
conservation residual reduces to the pure EOS-table term, with every
numerical channel at the few$\times10^{-5}$ level.  Re-metering the
ledger to $T^{\,n}$ instead would close the identity by construction
while measuring nothing, and is rejected for the same reason as for
$C$.

\section{Summary}

\begin{itemize}
\item The conservation residual is a checksum, not an error bar: balanced
errors pass it.  The exact budget $\dE = \epsA + \epsS + C$
(Eq.~\ref{eq:budget}) converts one confounded number into three
attributable ones, one per subsystem (EOS tables, transport, hydrostatic
scheme).
\item Ledger corrections (Section~\ref{sec:fixes}) made the accounting
internally consistent.  The honest residual is \emph{larger} than the
legacy number, because the legacy number benefited from bias cancellation.
\item $\epsA \simeq +9.3\times10^{-3}\,\Erad$ (homogeneous) /
$+10.0\times10^{-3}$ (inhomogeneous) is the EOS-table Maxwell defect:
resolution-independent, currently irreducible, and the floor to which the
residual converges.  It is also a concrete, measurable target for future
table construction.
\item $C = -4.7/N\,\Erad$ was the entire resolution-dependent error, and
is removed (200$\times$) by centering one pressure factor in the Henyey
gravity term.  That change simultaneously accelerates observable
convergence, to the point that $N=300$ matches the extrapolated continuum
limit.
\item Quote $|C|\simeq 4.7/N$ (legacy scheme) as the numerical energy
error, $\epsA + \epsS \simeq +1.1\%$ of $\Erad$ as the physics-input
systematic, and never tune $N$ to minimize the raw residual.
\end{itemize}

\end{document}
